% problem-000342 6 方解石与氟磷灰石晶胞
\section*{6. 方解石与氟磷灰石晶胞}
\textit{下列为分问。}

\subsection*{6-1}
⽂献中常⽤下图表达⽅解⽯的晶体结构：

（图：）

图中的平⾏六⾯体是不是⽅解⽯的⼀个晶胞？简述理由。

\subsection*{6-2}
⽂献中常⽤下图表达六⽅晶体氟磷灰⽯的晶体结构：

（图：）

该图是�轴投影图，位于图中⼼的球是氟，⼤球是钙，四⾯体是磷酸根（氧原⼦未画出）。试以此图为基础⽤粗线画出氟磷灰⽯晶胞的�轴投影图，设晶胞顶⻆为氟原⼦，其他原⼦可不补全。

\subsection*{6-3}
某晶体的晶胞参数为： $a = 2 5 0 . 4 \mathrm { p m }$ 1， $c = 6 6 6 . 1 \mathrm { p m }$ $\gamma = 1 2 0 ^ { \circ }$ ；原⼦A的原⼦坐标为 $( 0 , 0 , 1 / 2 )$ 和(1/3, 2/3, 0)，原⼦ B 的原⼦坐标为(1/3, 2/3, 1/2)和(0, 0, 0)。

\subsection*{6-3}
-1 试画出该晶体的晶胞透视图（设晶胞底⾯即��⾯垂直于纸⾯，A原⼦⽤ $^ { 6 6 } \bigcirc ^ { 5 }$ 表⽰，B原⼦⽤ $" \bullet \bullet \bullet \bullet$ 表⽰）。

\subsection*{6-3}
-2 计算上述晶体中A和B两原⼦间的最⼩核间距 $d ( \mathbf { A B } )$ o

\subsection*{6-3}
-3 共价晶体的导热是共价键的振动传递的。实验证实，该晶体垂直于�轴的导热性⽐平⾏于�轴的导热性⾼20倍。⽤上述计算结果说明该晶体的结构与导热性的关系。
